\section*{Chapitre \ref{ch:b2:plant-meristems} --- Le développement végétatif des plantes~: méristèmes et croissance}

\begin{solution}{exo:b2:plant-meristems:1}
Un \omterm{def:b2:plant-meristems:meristem}{méristème} est une population persistante de cellules indifférenciées
en division, dont dérivent les tissus de la plante. \omterm{def:b2:plant-meristems:meristem}{Méristème apical
caulinaire}~: feuilles, tige, \omterm{def:b2:angiosperm-reproduction:flower}{fleurs}~; \omterm{def:b2:plant-meristems:meristem}{méristème apical racinaire}~: la
racine~; \omterm{def:b2:plant-meristems:meristem}{cambium}~: xylème secondaire (\omterm{def:b2:plant-meristems:wood}{bois}) et phloème secondaire~;
phellogène~: le \omterm{def:b2:plant-meristems:wood}{périderme} (liège) de l'\omterm{def:b2:plant-meristems:wood}{écorce}.
\end{solution}

\begin{solution}{exo:b2:plant-meristems:2}
Coiffe (protection, lubrification, perception de la gravité, cellules
desquamées)~; \omterm{def:b2:plant-meristems:meristem}{méristème} avec le \omterm{prop:b2:plant-meristems:ram}{centre quiescent} (division)~; \omterm{prop:b2:plant-meristems:ram}{zone
d'élongation} (les cellules s'allongent de dix à vingt fois)~; zone de
\omterm{def:b2:cell-differentiation:differentiation}{différenciation} (\omterm{prop:b2:plant-meristems:ram}{poils absorbants}, xylème et phloème mûrs). Les racines
latérales naissent du péricycle de la racine mûre, en face des pôles du
xylème.
\end{solution}

\begin{solution}{exo:b2:plant-meristems:3}
Des plants de haricot furent décapités~: les bourgeons latéraux se
développèrent. Sur un deuxième lot, le moignon reçut un bloc de gélose
contenant de l'auxine~: les bourgeons restèrent dormants~; sur le lot
témoin, un bloc de gélose seule~: les bourgeons se développèrent.
Conclusion~: c'est l'auxine issue de l'apex qui inhibe les bourgeons
latéraux --- la \omterm{prop:b2:plant-meristems:apical}{dominance apicale} est hormonale.
\end{solution}

\begin{solution}{exo:b2:plant-meristems:4}
Soixante ans~; l'aubier correspond aux quinze derniers cernes~; sont
vivants le \omterm{def:b2:plant-meristems:meristem}{cambium}, le phloème, les cellules des rayons et du parenchyme
de l'aubier, et le phellogène --- une mince enveloppe autour d'un cœur
mort.
\end{solution}

\begin{solution}{exo:b2:plant-meristems:5}
$0.1\times(P - 0.3)$~: \qty{0.01}{h^{-1}} (doublement en
$\ln 2/0.01 =
\qty{69}{h}$), \qty{0.03}{h^{-1}} (\qty{23}{h}), \qty{0.05}{h^{-1}}
(\qty{14}{h}). À \qty{0.25}{MPa}, la turgescence est inférieure au
seuil~: la croissance s'arrête, avant tout flétrissement.
\end{solution}

\begin{solution}{exo:b2:plant-meristems:6}
Feuilles 1 à 8 à 0, 137.5, 275, 52.5, 190, 327.5, 105,
\qty{242.5}{\degree}. Dans l'ordre~: 0, 52.5, 105, 137.5, 190, 242.5,
275, 327.5 --- le plus petit écart est de \qty{32.5}{\degree}. Chaque
feuille se place dans un espace laissé par les précédentes~: les huit se
font aussi peu d'ombre que possible et la couronne remplit
uniformément le disque de lumière.
\end{solution}

\begin{solution}{exo:b2:plant-meristems:7}
$2\pi r h\Delta r$~: à \qty{5}{cm}, $2\pi\times 0.05\times 12\times
0.003 = \qty{0.011}{m^3}$, \qty{5.7}{kg}, \qty{2.8}{kg} de carbone~; à
\qty{40}{cm}, \qty{0.090}{m^3}, \qty{45}{kg}, \qty{23}{kg} de carbone
--- huit fois plus pour la même largeur de cerne.
\end{solution}

\begin{solution}{exo:b2:plant-meristems:8}
\emph{clv3}~: plus de frein --- \omterm{def:b2:plant-meristems:meristem}{méristèmes} énormes, feuilles et pièces
florales surnuméraires, tiges fasciées. \emph{wus}~: plus
d'accélérateur --- le \omterm{def:b2:plant-meristems:meristem}{méristème} s'épuise après quelques feuilles et se
reconstitue à plusieurs reprises. Double mutant~: le \omterm{def:b2:meiosis-heredity:vocabulary}{phénotype}
\emph{wus}, donc WUS agit en aval de CLV3~: le rôle de CLV3 est de
réprimer WUS, et sans WUS il n'y a rien à réprimer.
\end{solution}

\begin{solution}{exo:b2:plant-meristems:9}
Transport~: \qty{20}{h}. Diffusion~: $(0.2)^{2}/(2\times 10^{-9}) =
\qty{2e7}{s}$, environ 230 jours. Le transport polarisé rend un signal
hormonal utilisable sur toute la longueur d'une plante en moins d'une
journée, et lui donne une direction.
\end{solution}

\begin{solution}{exo:b2:plant-meristems:10}
\qty{3}{cm} par jour de tissu mûr, avec des cellules de
\qty{200}{\micro m}~: \num{150} cellules par file et par jour. Cent
cellules en division doivent donc se diviser chacune 1.5 fois par jour~:
un cycle d'environ \qty{16}{h}.
\end{solution}

\begin{solution}{exo:b2:plant-meristems:11}
Les gibbérellines allongent les entre-nœuds~; une plante incapable d'y
répondre a des entre-nœuds courts et une tige courte. Une forte
fertilisation azotée fait grandir un blé haut, qui s'alourdit jusqu'à
se coucher (verse) et pourrir~; une tige courte et rigide reste debout,
et place le carbone supplémentaire dans le grain plutôt que dans la
paille, si bien que l'indice de récolte augmente. Dans la nature, une
plante courte est dominée et ombragée par ses voisines, et elle perd.
\end{solution}

\begin{solution}{exo:b2:plant-meristems:12}
Les \omterm{def:b2:plant-meristems:meristem}{méristèmes} permettent à une plante d'ajouter des organes là et
quand les conditions les favorisent --- des feuilles vers la lumière,
des racines vers l'eau --- et d'abandonner ceux qui ne sont pas
rentables~; chaque bourgeon axillaire est un point de croissance de
rechange, si bien que la plante survit au broutage, au feu et à la
taille en repoussant~; et, comme les \omterm{def:b2:plant-meristems:meristem}{méristèmes} restent perpétuellement
embryonnaires, un \omterm{def:b2:asexual-reproduction:clone}{génet} n'a pas de durée de vie fixée. Le coût~: une
plante ne peut ni se déplacer ni réorganiser ce qu'elle a construit ---
un organe une fois placé reste où il est, et un corps construit par
accrétion doit porter avec lui son passé mort (le duramen).
\end{solution}

\begin{solution}{pb:b2:plant-meristems:1}
\textbf{1.} $150/3 = 50$ feuilles par apex~; \num{10000} pour l'arbre.
\textbf{2.} 0, 137.5, 275, 52.5 et \qty{190}{\degree}.
\textbf{3.} $8\times 137.5 = 1100 = 3\times 360 + 20$~; $13\times 137.5
= 1787.5 = 5\times 360 - 12.5$~; $21\times 137.5 = 2887.5 = 8\times 360
+ 7.5$~: la 22\textsuperscript{e} feuille est la première à moins de
\qty{10}{\degree} de la première.
\textbf{4.} Chaque nouvelle feuille se place dans le plus grand espace
libre, si bien que les feuilles se font le moins d'ombre possible~;
c'est l'auxine, dont le transport polarisé (transporteurs PIN) draine
le voisinage de chaque primordium, qui fixe l'espacement.
\textbf{5.} $\num{10000}\times 80\,\text{cm}^{2} = \qty{80}{m^2}$~: toute
la couronne est faite des feuilles de la saison.
\textbf{6.} La totalité~; les premières feuilles sont construites à
partir de l'amidon stocké dans le bois et les racines l'année
précédente.
\textbf{7.} Un apex triplé produit des primordiums plus vite --- deux ou
trois par plastochrone --- si bien que le nombre de feuilles et la
surface foliaire pourraient tripler, mais les pousses seraient
fasciées, avec des tiges en ruban et une phyllotaxie désordonnée, et une
grande partie de la surface supplémentaire se ferait de l'ombre à
elle-même.
\textbf{8.} Jour~: $0.05\times(0.55 - 0.35) = \qty{0.01}{h^{-1}}$~;
nuit~: $0.05\times 0.40 = \qty{0.02}{h^{-1}}$.
\textbf{9.} Jour~: $2e^{0.12} = \qty{2.25}{cm}$, soit \qty{0.25}{cm}
ajoutés~; nuit~: $2.25e^{0.24} = \qty{2.87}{cm}$, soit \qty{0.61}{cm}
ajoutés.
\textbf{10.} Vitesse moyenne \qty{0.015}{h^{-1}} pendant \qty{96}{h}~:
$2e^{1.44}
= \qty{8.4}{cm}$.
\textbf{11.} Jour~: $0.05\times 0.30 = \qty{0.015}{h^{-1}}$~; nuit~:
$0.05\times 0.50 = \qty{0.025}{h^{-1}}$.
\textbf{12.} $0.05\times 0.05 = \qty{0.0025}{h^{-1}}$, le quart de la
vitesse diurne normale. L'ajustement osmotique augmente la teneur de la
cellule en solutés, si bien que, au même potentiel hydrique, la
turgescence remonte au-dessus du seuil.
\textbf{13.} Le jour, la transpiration abaisse le potentiel hydrique de
la feuille et la turgescence avec lui, si bien que $P - Y$ est faible~;
la nuit, la turgescence est rétablie. Heure par heure, la croissance est
limitée par l'eau, non par le sucre.
\textbf{14.} $2\pi\times 0.10\times 8\times 0.005 = \qty{0.025}{m^3}$~;
\qty{12.6}{kg} de matière sèche.
\textbf{15.} \qty{6.3}{kg} de carbone.
\textbf{16.} $80\times 4\times 150 = \qty{48}{kg}$ de carbone.
\textbf{17.} $6.3/48 = \qty{13}{\%}$. Autres puits~: les feuilles
elles-mêmes, les racines, les branches et les rameaux, la respiration
de toutes les cellules vivantes (environ la moitié de ce qui est fixé),
les réserves, les \omterm{def:b2:angiosperm-reproduction:flower}{fleurs} et les \omterm{def:b2:angiosperm-reproduction:seed}{graines}.
\textbf{18.} Rayon \qty{0.20}{m}~: $2\pi\times 0.20\times 8\times
0.005 = \qty{0.050}{m^3}$, le double de cette année --- la circonférence
le long de laquelle travaille le \omterm{def:b2:plant-meristems:meristem}{cambium} a doublé (et l'arbre sera plus
haut).
\textbf{19.} Une année sèche ou froide~; ou une défoliation par des
insectes, le gel ou le feu. Une cause climatique rétrécit les mêmes
cernes chez tous les arbres de la région et se corrèle aux relevés
météorologiques~; un dommage touche un seul arbre ou un seul peuplement,
et le gel ou le feu laissent des cicatrices anatomiques dans le cerne.
\textbf{20.} Les bourgeons les plus proches de chaque coupe se
développent en quelques jours en plusieurs pousses~; au bout d'un mois,
l'arbre est plus buissonnant et plus dense, avec des pousses plus
nombreuses mais plus courtes.
\textbf{21.} L'auxine appliquée sur la coupe remplace le signal de
l'apex~: les bourgeons restent dormants et aucune ramification ne suit.
\textbf{22.} Chaque taille supprime les apex et libère les bourgeons
situés en dessous~; six fois par an, chaque libération ajoute un étage de
pousses courtes, et la surface s'en remplit.
\textbf{23.} Des bourgeons axillaires dormants de la souche et de
bourgeons adventifs formés par le \omterm{def:b2:plant-meristems:meristem}{cambium}~; chaque bourgeon est un
\omterm{def:b2:plant-meristems:meristem}{méristème} complet capable de reconstruire le système caulinaire, alors
que la tête d'un animal est un \omterm{def:b2:vertebrate-development:induction}{organisateur} unique et irremplaçable,
sans copie de réserve.
\textbf{24.} Davantage de cytokinines atteignant les bourgeons favorise
leur débourrement --- arroser une plante taillée la fait se ramifier
davantage.
\textbf{25.} 50 feuilles par apex, \num{10000} par arbre~; vitesses de
croissance de l'entre-nœud \qty{0.01}{h^{-1}} le jour et
\qty{0.02}{h^{-1}} la nuit~; \omterm{def:b2:plant-meristems:wood}{bois} de l'année \qty{0.025}{m^3},
\qty{12.6}{kg} de matière sèche, \qty{6.3}{kg} de carbone.
\end{solution}
